%% 
%% Copyright 2007-2024 Elsevier Ltd
%% 
%% This file is part of the 'Elsarticle Bundle'.
%% ---------------------------------------------
%% 
%% It may be distributed under the conditions of the LaTeX Project Public
%% License, either version 1.3 of this license or (at your option) any
%% later version.  The latest version of this license is in
%%    http://www.latex-project.org/lppl.txt
%% and version 1.3 or later is part of all distributions of LaTeX
%% version 1999/12/01 or later.
%% 
%% The list of all files belonging to the 'Elsarticle Bundle' is
%% given in the file `manifest.txt'.
%% 
%% Template article for Elsevier's document class `elsarticle'
%% with numbered style bibliographic references
%% SP 2008/03/01
%% $Id: elsarticle-template-num.tex 249 2024-04-06 10:51:24Z rishi $
%%
\documentclass[preprint,12pt]{elsarticle}

%% Use the option review to obtain double line spacing
%% \documentclass[authoryear,preprint,review,12pt]{elsarticle}

%% Use the options 1p,twocolumn; 3p; 3p,twocolumn; 5p; or 5p,twocolumn
%% for a journal layout:
%% \documentclass[final,1p,times]{elsarticle}
%% \documentclass[final,1p,times,twocolumn]{elsarticle}
%% \documentclass[final,3p,times]{elsarticle}
%% \documentclass[final,3p,times,twocolumn]{elsarticle}
%% \documentclass[final,5p,times]{elsarticle}
%% \documentclass[final,5p,times,twocolumn]{elsarticle}

%% For including figures, graphicx.sty has been loaded in
%% elsarticle.cls. If you prefer to use the old commands
%% please give \usepackage{epsfig}

%% The amssymb package provides various useful mathematical symbols
\usepackage{amssymb}
%% The amsmath package provides various useful equation environments.
\usepackage{amsmath}
%% The amsthm package provides extended theorem environments
%% \usepackage{amsthm}

%% The lineno packages adds line numbers. Start line numbering with
%% \begin{linenumbers}, end it with \end{linenumbers}. Or switch it on
%% for the whole article with \linenumbers.
%% \usepackage{lineno}

\journal{Computer Methods in Applied Mechanics and Engineering}

\begin{document}

\begin{frontmatter}

%% Title, authors and addresses

%% use the tnoteref command within \title for footnotes;
%% use the tnotetext command for theassociated footnote;
%% use the fnref command within \author or \affiliation for footnotes;
%% use the fntext command for theassociated footnote;
%% use the corref command within \author for corresponding author footnotes;
%% use the cortext command for theassociated footnote;
%% use the ead command for the email address,
%% and the form \ead[url] for the home page:
%% \title{Title\tnoteref{label1}}
%% \tnotetext[label1]{}
%% \author{Name\corref{cor1}\fnref{label2}}
%% \ead{email address}
%% \ead[url]{home page}
%% \fntext[label2]{}
%% \cortext[cor1]{}
%% \affiliation{organization={},
%%             addressline={},
%%             city={},
%%             postcode={},
%%             state={},
%%             country={}}
%% \fntext[label3]{}

\title{A New Asymptotic-Preserving Dual Formulation Finite-Volume Method for the Compressible Euler Equations}

%% use optional labels to link authors explicitly to addresses:
%% \author[label1,label2]{}
%% \affiliation[label1]{organization={},
%%             addressline={},
%%             city={},
%%             postcode={},
%%             state={},
%%             country={}}
%%
%% \affiliation[label2]{organization={},
%%             addressline={},
%%             city={},
%%             postcode={},
%%             state={},
%%             country={}}

%\author{} %% Author name
%
%%% Author affiliation
%\affiliation{organization={},%Department and Organization
%            addressline={}, 
%            city={},
%            postcode={}, 
%            state={},
%            country={}}


\author[NCSU]{Alina Chertock} %% Author name
\ead{chertock@math.ncsu.edu}



%% Author affiliation
\affiliation[NCSU]{organization={Department of Mathematics, North Carolina State University},%Department and Organization
	addressline={SAS Hall, 2108, 2311 Stinson Dr}, 
	city={Raleigh},
	postcode={27607}, 
	state={North Carolina},
	country={United States of America}}


\author[UMICH]{Smadar Karni} %% Author name
\ead{karni@umich.edu}

%% Author affiliation
\affiliation[UMICH]{organization={Department of Mathematics, University of Michigan},%Department and Organization
	addressline={2074 East Hall, 530 Church Street}, 
	city={Ann Arbor},
	postcode={48109}, 
	state={Michigan},
	country={United States of America}}

\author[SUSTECH]{Alexander Kurganov} %% Author name
\ead{alexander@sustech.edu.cn}

%% Author affiliation
\affiliation[SUSTECH]{organization={Department of Mathematics and Shenzhen International Center for Mathematics, Southern University of Science and Technology},%Department and Organization
	addressline={No. 1088 Xueyuan Avenue, Nanshan District}, 
	city={Shenzhen},
	postcode={518055}, 
	state={Guangdong},
	country={China}}


\author[NCSU]{Lorenzo Micalizzi\corref{corref}} %% Author name
\ead{lmicali@ncsu.edu}
\cortext[corref]{Corresponding author}

%% Author affiliation
%\affiliation[NCSU]{organization={Department of Mathematics, North Carolina State University},%Department and Organization
%	addressline={SAS Hall, 2108, 2311 Stinson Dr}, 
%	city={Raleigh},
%	postcode={27607}, 
%	state={North Carolina},
%	country={United States of America}}

%% Abstract
\begin{abstract}
The paper focuses on the development of numerical methods for the compressible Euler equations. It is well-known that if the Mach number is
small, the system becomes stiff and hence explicit schemes suffer from severe time-step restrictions, making them inefficient or even
impractical. Our objective is to develop an asymptotic preserving (AP) scheme that remains uniformly accurate and stable across all Mach
numbers.

Instead of the conservative hyperbolic flux splitting approach, which is widely used to design AP schemes, we consider a primitive
(nonconservative) formulation and introduce a nonconservative hyperbolic splitting. The resulting system is discretized using a
semi-implicit approach: the stiff part is handled semi-implicitly using second-order central differences, while the nonstiff part is treated
explicitly using a second-order path-conservative central-upwind discretization. A key feature of our method is that the pressure at each
time level is computed by solving a well-posed Poisson-type elliptic equation, thereby enforcing the AP property. Simultaneously, we evolve
the conservative form of the system using a semi-discrete central-upwind (CU) scheme. At the end of each stage of the time discretization,
we perform a special post-processing that selects the appropriate numerical solution depending on the Mach number. This guarantees that in
low-Mach-number regimes, the solution is obtained by the AP nonconservative scheme, while in higher-Mach-number regimes, a sharp and
physically relevant solution is computed by the conservative CU scheme.

Numerical experiments confirm that the proposed AP scheme achieves the expected second order of accuracy and that the time-step constraint
is independent of the Mach number, making it a robust and efficient alternative to conventional explicit methods.
\end{abstract}

%%%Graphical abstract
%\begin{graphicalabstract}
%%\includegraphics{grabs}
%\end{graphicalabstract}
%
%%%Research highlights
%\begin{highlights}
%\item Research highlight 1
%\item Research highlight 2
%\end{highlights}

%% Keywords
\begin{keyword}
%% keywords here, in the form: keyword \sep keyword
Compressible Euler equations \sep low Mach number \sep asymptotic preserving (AP) schemes \sep hyperbolic splitting \sep semi-implicit
methods \sep deferred correction.


%% PACS codes here, in the form: \PACS code \sep code
\PACS 47.11.Df \sep 47.10.A- \sep 47.40.Hg \sep 02.60.Jh

%% MSC codes here, in the form: \MSC code \sep code
%% or \MSC[2008] code \sep code (2000 is the default)
\MSC 65M08 \sep 65M20 \sep 76M12 \sep 35L65 \sep 76N15 \sep 35B40

\end{keyword}

\end{frontmatter}

%% Add \usepackage{lineno} before \begin{document} and uncomment 
%% following line to enable line numbers
%% \linenumbers

%% main text
%%

%% Use \section commands to start a section
\section{Introduction}
%% If you have bib database file and want bibtex to generate the
%% bibitems, please use
%%
\bibliographystyle{elsarticle-num} 
\bibliography{ref}

%% else use the following coding to input the bibitems directly in the
%% TeX file.

%% Refer following link for more details about bibliography and citations.
%% https://en.wikibooks.org/wiki/LaTeX/Bibliography_Management

%\begin{thebibliography}{00}
%
%%% For numbered reference style
%%% \bibitem{label}
%%% Text of bibliographic item
%
%\bibitem{lamport94}
%  Leslie Lamport,
%  \textit{\LaTeX: a document preparation system},
%  Addison Wesley, Massachusetts,
%  2nd edition,
%  1994.
%
%\end{thebibliography}
\end{document}

\endinput
%%
%% End of file `elsarticle-template-num.tex'.
